% Source-derived chapter 10: Comparison of two cycle classes
% Original source: higher-siegel-weil-unitary-nonsingular-terms
\chapter{Comparison of two cycle classes}
\label{higher-siegel-weil-unitary-nonsingular-terms-chapter-10}
\source{higher-siegel-weil-unitary-nonsingular-terms}

\begin{remark}[Source section]
\label{higher-siegel-weil-unitary-nonsingular-terms-chapter-10-source}
\source{higher-siegel-weil-unitary-nonsingular-terms}
\source{higher-siegel-weil-unitary-nonsingular-terms}
This chapter reproduces the corresponding section of the retrieved
LaTeX source, including its definitions, statements, examples, and
proofs. It has not yet been translated into Lean.
\notready
\end{remark}

% The source section title was: Comparison of two cycle classes
\label{s:compare cycles}
\source{higher-siegel-weil-unitary-nonsingular-terms}
The goal in this section is to show the following theorem.

\begin{thm}
\source{higher-siegel-weil-unitary-nonsingular-terms}
\notready\label{th:compare KR} Let $\cE\in \Bun_{\GL_{n}'}(k)$ and $a\in \cA_{\cE}(k)$. Let $s:\cE'=\op_{i=1}^{n}\cL_{i}\inj \cE$ be a good framing of $(\cE,a)$ in the sense of Definition \ref{def: auxiliary E'}. Let $a':\cE'\to\s^{*}(\cE')^{\vee}$ be the Hermitian map induced from $a$. Then we have an equality in the Chow group
\source{higher-siegel-weil-unitary-nonsingular-terms}
\begin{equation}\label{KR vs ShtM} 
\source{higher-siegel-weil-unitary-nonsingular-terms}
\z^{r}_{\cL_{1},\cdots, \cL_{n}}(a')|_{\cZ^{r}_{\cE}(a)}=[\Sht^{r}_{\cM(n,n)}]|_{\cZ^{r}_{\cE}(a)}\in \Ch_{0}(\cZ^{r}_{\cE}(a)).
\end{equation}
Here the restriction on the RHS is via \eqref{eq: hitchin shtuka decomposition E and A inj}, noting that $\cZ_{\cE}^r(a)^{\circ}=\cZ_{\cE}^r(a)$, and the class $[\Sht^{r}_{\cM(n,n)}]$ is as in Definition \ref{def: hitchin shtukas class}.% plus the observation that $a \in \cA_{\cE}(k)$ implies $\cZ^{r}_{\cE}(a) = \cZ^{r}_{\cE}(a)^\circ$ (cf. Definition \ref{def: A(k)} and Definition \ref{def:Z(0)}).  

In particular, the cycle class $[\cZ^{r}_{\cE}(a)]$ as in Definition-Proposition \ref{defn:general KR} is well-defined (i.e., independent of the choice of a good framing). 
\end{thm}

Below we consider the case where $X'$ is geometrically connected. At the end of this section (\S\ref{ss:compare cycle split}) we comment on how to modify the argument in the case $X'=X\coprod X$ or $X'=X_{k'}$, where $k'/k$ is the quadratic extension.

\subsection{First reductions}
For a vector bundle $\cE$ on $X'$ let $\mu_{\min}(\cE)\in \Q$ be the smallest slope that appears in the Harder-Narasimhan filtration of $\cE$.  For $\cE$ of rank $n$ and $a\in \cA_{\cE}(k)$, a good framing $s:\op_{i=1}^{n}\cL_{i}\inj \cE$ for $(\cE,a)$ is called {\em very good} if it satisfies the additional condition
\begin{enumerate}
\item[(3)] $\mu_{\min}(\cE)>\max\{\deg\cL_{i}+2g'-1\}_{1\le i\le n}$.
\end{enumerate}
%\mnote{Y: if $X'=X\coprod X$,  the inequality should be understood as two inequalities, one on each component of $X'$ and $g'=g$.}

Most of the work in this section will be devoted to proving the slightly weaker statement below. 
\begin{thm}
\source{higher-siegel-weil-unitary-nonsingular-terms}
\notready\label{th:compare KR weak} Suppose $X'$ is connected. Then the identity \eqref{KR vs ShtM} holds if $s:\cE'=\op_{i=1}^{n}\cL_{i}\inj \cE$ is a very good framing of $(\cE,a)$.
\source{higher-siegel-weil-unitary-nonsingular-terms}
\end{thm}


\begin{lemma}
\source{higher-siegel-weil-unitary-nonsingular-terms}
\notready\label{l:weak implies strong} Theorem \ref{th:compare KR weak} implies Theorem \ref{th:compare KR}.
\source{higher-siegel-weil-unitary-nonsingular-terms}
\end{lemma}
\begin{proof} Choose effective divisors $D_{i}$ on $X'$ ($1\le i\le n$) such that $\nu(D_{1}+\cdots+D_{n})$ is multiplicity-free and disjoint from $\nu^{-1}(D_{a})$. Let $\cL'_{i}=\cL_{i}(-D_{i})$. When the $D_{i}$'s have sufficiently large degree,  the resulting map
\begin{equation}
s': \op_{i=1}^{n}\cL'_{i}\inj \op_{i=1}^{n}\cL_{i}\inj \cE
\end{equation}
is a very good framing. Let $a''$ be the induced Hermitian map $\op\cL'_{i} \rightarrow \sigma^* (\op\cL'_{i})^{\vee}$. By Theorem \ref{th:compare KR weak} we have
\begin{equation}
\z^{r}_{\cL'_{1},\cdots, \cL'_{n}}(a'')|_{\cZ^{r}_{\cE}(a)}=[\Sht^{r}_{\cM(n,n)}]|_{\cZ^{r}_{\cE}(a)}.
\end{equation}
Therefore, to prove \eqref{KR vs ShtM} it suffices to show 
\begin{equation}
\z^{r}_{\cL_{1},\cdots, \cL_{n}}(a')|_{\cZ^{r}_{\cE}(a)}=\z^{r}_{\cL'_{1},\cdots, \cL'_{n}}(a'')|_{\cZ^{r}_{\cE}(a)}\in\Ch_{0}(\cZ^{r}_{\cE}(a)).
\end{equation}
Let $U'$ be the complement of $\cup_{i=1}^{n}\supp(D_{i}+\s D_{i})$ in $X'$. By construction, $U'$ contains $\nu^{-1}(D_{a})$, therefore $\cZ_{\cE}^{r}(a)|_{U'^{r}}=\cZ_{\cE}^{r}(a)$ by Lemma \ref{l:KR cycle legs}. Let $\z_{i}=[\cZ^{r}_{\cL_{i}}(a'_{ii})^{\circ}]|_{U'^{r}}\in \Ch_{r(n-1)}(\cZ^{r}_{\cL_{i}}(a'_{ii})^{\circ}|_{U'^{r}})$. Similarly define $\z'_{i}$ using $(\cL'_{i}, a''_{ii})$. Then it suffices to show the equality
\begin{equation}\label{zz}
\source{higher-siegel-weil-unitary-nonsingular-terms}
(\z_{1}\cdot \z_{2}\cdot \dots \cdot\z_{n})|_{\cZ^{r}_{\cE}(a)}=(\z'_{1}\cdot \z'_{2}\cdot \dots \cdot\z'_{n})|_{\cZ^{r}_{\cE}(a)}\in \Ch_{0}(\cZ^{r}_{\cE}(a)),
\end{equation}
where the intersection products are taken over $\Sht^{r}_{U(n)}|_{U'^{r}}$. Applying Lemma \ref{l:change E}  to each injection $\cL'_{i}\inj \cL_{i}$, we see that $\cZ^{r}_{\cL_{i}}(a'_{ii})|_{U'^{r}}\inj \cZ^{r}_{\cL'_{i}}(a''_{ii})|_{U'^{r}}$ is open and closed. Therefore $\cZ^{r}_{\cL_{i}}(a'_{ii})^{\circ}|_{U'^{r}}\inj \cZ^{r}_{\cL'_{i}}(a''_{ii})^{\circ}|_{U'^{r}}$ is open. This shows that the fundamental class $\z_{i}$ is the open restriction of $\z'_{i}$ to $\cZ^{r}_{\cL_{i}}(a'_{ii})^{\circ}|_{U'^{r}}$. The equality \eqref{zz} then follows.
\end{proof}



\subsection{Auxiliary moduli spaces}
Let $\un d=(d_{i})_{1\le i\le n}\in \Z^{n}_{\ge0}$ and $e\in \Z_{\ge0}$. Write $d=\sum d_{i}$. 

Recall that $\cM_{e}\subset \cM(n,n)$ is the open-closed substack where $\chi(X',\cE)=-e$. Let $\cM_{\un d}$ be the moduli stack classifying $(\{\cL_{i}\}_{1\le i\le n}, (\cF,h),\{t'_{i}:\cL_{i}\to \cF\}_{1\le i\le n})$ where
\begin{itemize}
\item $\cL_{i}$ is a line bunde on $X'$ with $\chi(X',\cL_{i})=-d_{i}$ for $1\le i\le n$;
%\item The composition $a'_{ii}: \cL_{i}\xr{t'_{i}}\cF\xr{h}\s^{*}\cF^{\vee}\xr{\s^{*}t'^{\vee}_{i}}\s^{*}\cL_{i}^{\vee}$ is nonzero for $1\le i\le n$;
\item $(\cF,h)\in \Bun_{U(n)}$ satisfying
\begin{equation}\label{slope F}
\source{higher-siegel-weil-unitary-nonsingular-terms}
\mu_{\min}(\cF)>\max\{-d_{i}+3g'-2\}_{1\le i\le n}.
\end{equation}
\item For each $1\le i\le n$, $t'_{i}: \cL_{i}\to \cF$ is an injective map (fiberwise over the test scheme).
\end{itemize}
We define $\Hk^{r}_{\cM_{\un d}}$ to be the moduli stack classifying
\begin{equation*}
(\{\cL_{i}\}_{1\le i\le n}, (\{x'_{i}\}_{1\le i\le r}, \{(\cF_{j},h_{j})\}_{0\le j\le r})\in \Hk^{r}_{U(n)},\{t'_{ij}:\cL_{i}\to \cF_{j}\})
\end{equation*}
where  $\cL_{i}$ are the same  as in $\cM_{\un d}$, each $\cF_{j}$ satisfies the analogue of \eqref{slope F} with $\cF$ replaced by $\cF_{j}$, and fiberwise injective maps $t'_{ij}:\cL_{i}\to \cF_{j}$ are compatible with the isomorphisms between $\cF_{j-1}$ and $\cF_{j}$ away from $x'_{i}$.

\begin{lemma}
\source{higher-siegel-weil-unitary-nonsingular-terms}
\notready\label{l:low deg M sm} The stacks $\cM_{\un d}$ and $\Hk^{1}_{\cM_{\un d}}$ are smooth stacks of pure dimension $dn - (n^2-2n)(g-1)$.
\source{higher-siegel-weil-unitary-nonsingular-terms}
\end{lemma}
\begin{proof}
We first prove the statement for $\cM_{\un d}$. 
%The automorphisms of an $S$-point $(t':\op\cL_{i}\inj \cF,h)\in \cM_{\un d}$ is contained in $\prod_{i=1}^{n}\Aut(\cL_{i}, a'_{ii})\cong(\mu_{2})^{n}$. This shows that $\cM_{\un d}$ is Deligne-Mumford.  
Consider the map $\g_{\un d}: \cM_{\un d}\to \prod_{i=1}^{n}\Pic^{-d_{i}+g'-1}_{X'}\times \Bun_{U(n)}$ sending $(\{\cL_{i}\},(\cF,h),\{t_i\})$ to $(\{\cL_{i}\},(\cF,h))$.  For $(\cF,h)\in \Bun_{U(n)}$ and $\cL_{i}\in \Pic^{-d_{i}+g'-1}_{X'}$, the condition $\mu_{\min}(\cF)\ge\max\{-d_{i}+3g'-2\}_{1\le i\le n}=\max\{\deg\cL_{i}+2g'-1\}_{1\le i\le n}$ guarantees that $\Ext^{1}(\cL_{i}, \cF)=\Hom(\cF,\cL_{i}\ot\om_{X'})^{\vee}=0$. Noting that $\deg \cF = n(g'-1)$, the Riemann-Roch formula implies that $\g_{\un d}$ exhibits $\cM_{\un d}$ as an open substack of a vector bundle of rank $\dim\Hom(\op\cL_{i}, \cF)=-n\sum_{i}\deg\cL_{i}=dn-n^{2}(g'-1)$ over the base. In particular, $\cM_{\un d}$ is smooth and equidimensional. Since $\dim \Pic^{-d_{i}+g'-1}_{X'}=g'-1$ and $\dim \Bun_{U(n)}=n^{2}(g-1)$, we conclude that 
\begin{align*}
\dim\cM_{\un d} & =dn-n^{2}(g'-1) + n (g'-1) + n^{2}(g-1) = dn- (n^2-2n)(g-1).
\end{align*}

The argument for $\Hk^{1}_{\cM_{\un d}}$ is similar. The natural map 
$$
\Hk^{1}_{\cM_{\un d}}\to \prod_{i=1}^{n}\Pic^{-d_{i}+g'-1}_{X'}\times\Hk^{1}_{U(n)}
$$
exhibits $\Hk^{1}_{\cM_{\un d}}$ as an open substack of a vector bundle of rank $\dim\Hom(\op\cL_{i}, \cF_{0}\cap \cF_{1})=dn-n$ (using that $\deg(\cF_{0}\cap\cF_{1})=n(g'-1)-1$) over the base. Here we need the stronger inequality $\mu_{\min}(\cF_{0})>\max\{-d_{i}+3g'-2\}$ to guarantee $\mu_{\min}(\cF_{0}\cap\cF_{1})\ge\max\{-d_{i}+3g'-2\}$. In particular, $\Hk^{1}_{\cM_{\un d}}$ is smooth and equidimensional, and 
\begin{align*}
\dim \Hk^{1}_{\cM_{\un d}} & =dn-n -n^{2}(g'-1)+ n(g'-1)+\dim \Hk^{1}_{U(n)}  \\
&= dn-n -2n^{2}(g-1)+ 2n(g-1) + n + n^2(g-1) =\dim\cM_{\un d},
\end{align*}
as desired. 
\end{proof}

Let $\cM_{\un d,e}$ be the  moduli stack of $(\{\cL_{i}\}_{1\le i\le n}, \cE, \cF,h,s,t)$ where 
\begin{itemize}
\item $\cL_{i}\in \Pic_{X'}$ satisfies $\chi(X',\cL_{i})=-d_{i}$ for $i=1,\cdots, n$; 
\item $\cE\in \Bun_{\GL_n'}$ satisfies $\chi(X',\cE)=-e$;
\item $(\cF,h)\in \Bun_{U(n)}$; 
\item $t: \cE\to \cF$ is an injective map;
\item $s: \op_{i=1}^{n}\cL_{i}\to\cE$ is a very good framing for $(\cE,a)$, where $a=\s^{*}t^{\vee}\c h\c t$ is the induced Hermitian map on $\cE$.
\end{itemize}
Note that being a very good framing requires $-d_{i}<\mu_{\min}(\cE)-(3g'-2)$ for all $i$, which imposes an open condition on $\cE$.  We view $\cM_{\un d,e}$ as a correspondence
\begin{equation}\label{corr M}
\source{higher-siegel-weil-unitary-nonsingular-terms}
\xymatrix{ & \cM_{\un d, e}\ar[dl]_{w_{1}}\ar[dr]^{w_{2}}\\
\cM_{\un d} & &\cM_{e}\,.}
\end{equation}
Here $w_{1}$ records $\op_{i=1}^{n}\cL_{i}\xr{t\circ s}\cF$ and $w_{2}$ records $\cE\xr{t}\cF$.

We denote the Hitchin bases for $\cM_{\un d}, \cM_{e}$ and $\cM_{\un d,e}$ by $\cA_{\un d}$,  $\cA_{e}$ and $\cA_{\un d,e}$ respectively. Here $\cA_{\un d}$ parametrizes $(\{\cL_{i}\}_{1\le i\le n},a'=(a'_{ij}))$ (where $a'_{ij}: \cL_{j}\to \s^{*}\cL_{i}^{\vee}$) such that $a':\op\cL_{i}\to \s^{*}(\op\cL_{i})^{\vee}$ is an injective Hermitian map.  The base $\cA_{e}$ classifies $(\cE, a)$ with $a$ an injective Hermitian map. The base $\cA_{\un d,e}$ is the moduli stack of $(\{\cL_{i}\}_{1\le i\le n}, \cE, s,a)$ where  $(\cE,a)\in \cA_{e}$, $s: \op_{i=1}^{n}\cL_{i}\to\cE$ is a very good framing of $(\cE,a)$ and $\cL_{i}\in \Pic_{X'}$ with $\chi(X',\cL_{i})=-d_{i}$.  We view $\cA_{\un d,e}$ as a correspondence 
\begin{equation}\label{corr A}
\source{higher-siegel-weil-unitary-nonsingular-terms}
\xymatrix{ & \cA_{\un d, e}\ar[dl]_{u_{1}}\ar[dr]^{u_{2}}\\
\cA_{\un d} & &\cA_{e}\,.}
\end{equation}
We have Hitchin maps
\begin{eqnarray*}
f_{\un d}: &\cM_{\un d}\to \cA_{\un d},\\
f_{e}: &\cM_{e}\to \cA_{e},\\
f_{\un d, e}: &\cM_{\un d, e}\to \cA_{\un d,e}.
\end{eqnarray*} 
These maps together give a map of correspondences \eqref{corr M} to \eqref{corr A}. Note $f_{\un d}$ is not necessarily proper because we have imposed an open condition on the minimal slope of $\cF$.

Similarly we define the Hecke version $\Hk^{r}_{\cM_{\un d ,e}}$ of $\cM_{\un d,e}$ as the moduli stack of $(\{x'_{i}\}_{1\le i\le r}, \{\cE\to\cF_{i}\}_{0\le i\le r})\in \Hk^{r}_{\cM_{e}}$ together with a very good framing $s:\op \cL_{i}\inj \cE$ for $(\cE,a)$ with $\chi(X',\cL_{i})=-d_{i}$. Again we view $\Hk^{r}_{\cM_{\un d,e}}$ as a correspondence
\begin{equation*}
\xymatrix{ & \Hk^{r}_{\cM_{\un d,e}}\ar[dl]_{h_{1}}\ar[dr]^{h_{2}}\\
\Hk^{r}_{\cM_{\un d}} & &\Hk^{r}_{\cM_{e}}\,.}
\end{equation*}


\begin{lemma}
\source{higher-siegel-weil-unitary-nonsingular-terms}
\notready\label{l:uwh et}
\source{higher-siegel-weil-unitary-nonsingular-terms}
The maps $w_{1},u_1$ and $h_{1}$ are \'etale.
\end{lemma}
\begin{proof}
We first prove that $u_{1}$ is \'etale.  Let $(X_{d-e}\times X_{e})^{\hs}$ be the open subscheme of divisors $(D_{1},D_{2})\in X_{d-e}\times X_{e}$ such that $D_{1}$ is multiplicity-free and disjoint from $D_{2}$; let $(X'_{d-e}\times X_{e})^{\hs}$ be the preimage of $(X_{d-e}\times X_{e})^{\hs}$ in $X'_{d-e}\times X_{e}$. We have a map $\a: X'_{d-e}\times X_{e}\to X_{d}$ sending $(D'_{1}, D_{2})$ to $\nu(D_{1})+D_{2}$. Let $\a^{\hs}$ be the restriction of $\a$ to  $(X'_{d-e}\times X_{e})^{\hs}$. By factorizing $\a^{\hs}$ as the composition 
\begin{equation}
 (X'_{d-e}\times X_{e})^{\hs} \xr{\nu_{d-e}\times \Id}  (X_{d-e}\times X_{e})^{\hs}\xr{\mrm{add}}X_{d} 
 \end{equation}
we see that $\a^{\hs}$ is \'etale.  From the definition we have a map
\begin{equation}\label{Ade}
\source{higher-siegel-weil-unitary-nonsingular-terms}
j=(u_{1}, j'_{d-e}, j_{e}): \cA_{\un d,e}\to \cA_{\un d}\times_{X_{d}}(X'_{d-e}\times X_{e})^{\hs}.
\end{equation}
where $j'_{d-e}: \cA_{\un d,e}\to X'_{d-e}$ sends $(\{\cL_{i}\}, \op \cL_{i}\xr{s}\cE,a)$ to $\Div(s)$ (the divisor of $\det(s)$) and $j_{e}: \cA_{\un d,e}\to X_{e}$ sends it to $D_{a}$  (see Definition \ref{def:Da}).  The map $\cA_{\un d}\to X_{d}$ used in the fiber product records the divisor $D_{a'}$ of the Hermitian map $a'$ on $\op \cL_{i}$. We claim that $j$ is an open immersion. Indeed, given $(\{\cL_{i}\},a')\in \cA_{\un d}$ and $(D'_{1}, D_{2})\in (X'_{d-e}\times X_{e})^{\hs}$ such that $\nu(D_{1})+D_{2}=D_{a'}$, by the disjointness of $D'_{1},\s(D'_{1})$ and $\nu^{-1}(D_{2})$, there is one and only one coherent sheaf $\cE$ such that $\op\cL_{i}\subset \cE\subset \s^{*}(\op\cL_{i})^{\vee}$, $\cE/\op\cL_{i}$ is supported on $D'_{1}$, and $\s^{*}(\op\cL_{i})^{\vee}/\cE$ is supported away from $D'_{1}$. This would give a very good framing of $\cE$ if the open condition $\mu_{\min}(\cE)>\max\{-d_{i}+3g'-2\}_{1\le i\le n}$ is satisfied. This shows that $j$ is an open immersion. Since $\a^{\hs}$ is \'etale, we conclude that $\cA_{\un d,e}$ is \'etale over $\cA_{\un d}$.

To show $w_{1}$ is \'etale, we observe that $\cM_{\un d,e}\cong \cM_{\un d}\times_{\cA_{\un d}}\cA_{\un d,e}$. Since $u_{1}$ is \'etale, so is $w_{1}$. 

%we consider the similarly defined map
%\begin{equation}
%j': \cM_{\un d,e}\to \cM_{\un d}\times_{X_{d}}(X'_{d-e}\times X_{e})^{\hs}
%\end{equation}
%and it is easy to show that $j'$ is also an open immersion. Hence  $\cM_{\un d,e}$ is \'etale over $\cM_{\un d}$.

%Consider the diagram
%\begin{equation}
%\xymatrix{\cN_{\un d, e} \ar[r] & \ar[r]^{\d}\ar[d] & (X'_{d-e}\times X_{e})^{\hs}\ar[r]^{\pr_{1}}\ar[d] & (X'_{d-e})^{\nu\c}\ar[d]^{\nu^{\c}_{d-e}}\\
%& \cN^{\flat}_{\un d,e}\ar[r]\ar[d] &  (X_{d-e}\times X_{e})^{\hs} \ar[d]^{\textup{add}}\ar[r]^{\pr_{1}} & (X_{d-e})^{\c}\\ 
%\cM_{\un d}\ar[r]^{\d'} & X_{d}}
%\end{equation}
%Here the map ``add'' is adding two divisors on $X$;  ; $X_{d-e}^{\c}$ is the multiplicity-free locus and $(X'_{d-e})^{\nu\c}$ is its preimage in $X'_{d-e}$. The map $\d'$ record the divisor $D_{a'}$ of the Hermitian form $a'$ on $\op \cL_{i}$ (see Definition \ref{def:Da}); $\d$ records the divisor of $s$ and $D_{a}$ of the Hermitian form $a$ on $\cE$. The stack $ \cN^{\flat}_{\un d,e}$ is defined to make the lower square Cartesian. By the definition of a very good framing, the map $\d$ lands in $(X'_{d-e}\times X_{e})^{\hs}$ and the outer square is Cartesian. It is easy to see that both ``add'' and $\nu_{d-e}^{\c}$ are \'etale, therefore so is the composition of the left column, which is $w_{1}$.

Finally,  $\Hk^{r}_{\cM_{\un d,e}}$ is the open substack of $\Hk^{r}_{\cM_{\un d}}\times_{\cA_{\un d}}\cA_{\un d,e}$  where the legs avoid $\Div(s)$. Since $u_{1}$ is \'etale,  so is $h_{1}$.
\end{proof}


%\begin{lemma}
%The maps $f_{\un d}:\cM_{\un d}\to \cA_{\un d}$ and $f_{\un d,e}: \cM_{\un d,e}\to \cA_{\un d,e}$ are small. 
%\end{lemma}
%\begin{proof}
%We have a commutative diagram
%\begin{equation}
%\xymatrix{\cM_{\un d}\ar[d]^{f_{\un d}} &  \cM_{\un d,e}\ar[l]_{w_{1}}\ar[r]^{w_{2}}\ar[d]^{f_{\un d,e}}\ar[r] & \cM_{e}\ar[d]^{f_{e}}\\
%\cA_{\un d} & \cA_{\un d,e}\ar[l]_{u_{1}}\ar[r]^{u_{2}} & \cA_{e}
%}
%\end{equation}
%The right square is Cartesian by definition.  The map $\cA_{\un d,e}\to \cA_{e}$ is smooth (it is open in a vector bundle over $\cA_{e}$). Therefore the smallness of $f_{\un d,e}$ follows from that of $f_{e}$ proved in Corollary \ref{c:small}. The square on the left is not Cartesian, but both $u_{1}$ and $w_{1}$ are \'etale  by Lemma \ref{l:uwh et}. Since $f_{\un d,e}$ is  just shown to be small, from which we conclude that $f_{\un d}$ is also small.
%\end{proof}



\subsection{Auxiliary Hitchin shtukas}
We define $\Sht^{r}_{\cM_{\un d}}$ and $\Sht^{r}_{\cM_{\un d,e}}$ as the fiber product
\begin{equation}\label{defn ShtMde}
\source{higher-siegel-weil-unitary-nonsingular-terms}
\xymatrix{ \Sht^{r}_{\cM_{\un d}}\ar[d] \ar[r] & \Hk^{r}_{\cM_{\un d}}\ar[d]^{(\pr_{0},\pr_{r})}& \Sht^{r}_{\cM_{\un d,e}}\ar[d] \ar[r] & \Hk^{r}_{\cM_{\un d,e}}\ar[d]^{(\pr_{0},\pr_{r})}\\
\cM_{\un d}\ar[r]^-{(\Id, \Frob)} & \cM_{\un d}\times \cM_{\un d} & \cM_{\un d,e}\ar[r]^-{(\Id, \Frob)} & \cM_{\un d,e}\times \cM_{\un d,e}}
\end{equation}

The maps $w_{i}$ and $h_{i}$ induce maps
\begin{equation}\label{corr ShtM}
\source{higher-siegel-weil-unitary-nonsingular-terms}
\xymatrix{& \Sht^{r}_{\cM_{\un d,e}}\ar[dl]_{u_{1}}\ar[dr]^{u_{2}}\\
\Sht^{r}_{\cM_{\un d}} & & \Sht^{r}_{\cM_{e}}}
\end{equation}
The stack $\Sht^{r}_{\cM_{\un d,e}}$ decomposes into a disjoint union of open-closed substacks indexed by $(\{\cL_{i}\}, \op\cL_{i}\xr{s} \cE, a)\in \cA_{\un d,e}(k)$ 
\begin{equation}
\Sht^{r}_{\cM_{\un d,e}}=\coprod_{(\op\cL_{i}\to \cE, a)\in \cA_{\un d,e}(k)}\cZ^{r}_{\cE}(a).
\end{equation}
Correspondingly, the diagram \eqref{corr ShtM} decomposes into the disjoint union indexed by $\cA_{\un d,e}(k)$ of diagrams of the form
\begin{equation}
\xymatrix{ & \cZ^{r}_{\cE}(a)\ar@{^{(}->}[dl]_{u_{1}}\ar@{=}[dr]^{u_{2}}\\
\cZ^{r}_{\op\cL_{i}}(a')^{\hs} & & \cZ^{r}_{\cE}(a) }
\end{equation}
Here $\cZ^{r}_{\op\cL_{i}}(a')^{\hs}\subset \cZ^{r}_{\op\cL_{i}}(a')$ (where $a'$ is the Hermitian map on $\op\cL_{i}$ induced from $a$) is cut by the open condition $\mu_{\min}(\cF_{j})>\max\{-d_{i}+3g'-2\}$ for all $0\le j\le r$. From this description  and Corollary \ref{c:ZE open closed}, we see that:

\begin{lemma}
\source{higher-siegel-weil-unitary-nonsingular-terms}
\notready\label{l:ShtM isom}
\source{higher-siegel-weil-unitary-nonsingular-terms}
The map $u_{1}$ (resp. $u_{2}$), when restricted to each connected component of $\Sht^{r}_{\cM_{\un d,e}}$, is an isomorphism onto a connected component of $\Sht^{r}_{\cM_{\un d}}$ (resp. $\Sht^{r}_{\cM_{e}}$).
\end{lemma}


\subsection{Zero cycles on auxiliary Hitchin shtukas}
Similar to the definition of $[\Sht^{r}_{\cM_{e}}]$ given in \S\ref{ss:ShtM cycle}, we define $0$-cycles supported on $\Sht^{r}_{\cM_{\un d,e}}$ and $\Sht^{r}_{\cM_{\un d}}$ as follows.

We rewrite $\Sht^{r}_{\cM_{\un d}}$ as the fiber product
\begin{equation}\label{rewrite ShtM und}
\source{higher-siegel-weil-unitary-nonsingular-terms}
\xymatrix{\Sht^{r}_{\cM_{\un d}}\ar[d]\ar[r] & (\Hk^{1}_{\cM_{\un d}})^{r}\ar[d]^{(\pr_{0},\pr_{1})^{r} \times \Delta} \times \cM_{\un d}\\
(\cM_{\un d})^{r+1}\ar[r]^{\Phi_{\cM_{\un d}}} & (\cM_{\un d})^{2r+2}
}
\end{equation}
Here $\Phi_{\cM_{\un d}}=\Phi^{r}_{\cM_{\un d}}$ and the vertical maps are defined as in Definition \ref{def:Phi}. By the smoothness of $\Hk^{1}_{\cM_{\un d}}$ and $\cM_{\un d}$ proved in Lemma \ref{l:low deg M sm} and the dimension calculation there, we define
\begin{equation}
[\Sht^{r}_{\cM_{\un d}}]:=\Phi_{\cM_{\un d}}^{!}[(\Hk^{1}_{\cM_{\un d}})^{r} \times \cM_{\un d}]\in\Ch_{0}(\Sht^{r}_{\cM_{\un d}}).
\end{equation}

Similarly, using the Cartesian diagram
\begin{equation}\label{rewrite ShtMde}
\source{higher-siegel-weil-unitary-nonsingular-terms}
\xymatrix{\Sht^{r}_{\cM_{\un d, e}}\ar[d]\ar[r] & (\Hk^{1}_{\cM_{\un d, e}})^{r}\ar[d]^{(\pr_{0},\pr_{1})^{r} \times \Delta} \times \cM_{\un d, e}\\
(\cM_{\un d ,e})^{r+1}\ar[r]^{\Phi_{\cM_{\un d,e}}} & (\cM_{\un d, e})^{2r+2}
}
\end{equation}
and the smoothness and dimension calculations of $\Hk^{1}_{\cM_{\un d,e}}$ and $\cM_{\un d, e}$ (which follow from Lemma \ref{l:uwh et} and Lemma \ref{l:low deg M sm}), we define
\begin{equation}
[\Sht^{r}_{\cM_{\un d, e}}]:=\Phi_{\cM_{\un d,e}}^{!}[(\Hk^{1}_{\cM_{\un d, e}})^{r} \times \cM_{\un d, e}]\in\Ch_{0}(\Sht^{r}_{\cM_{\un d ,e}}).
\end{equation}


\begin{lemma}
\source{higher-siegel-weil-unitary-nonsingular-terms}
\notready\label{l:u1}
\source{higher-siegel-weil-unitary-nonsingular-terms}
We have $u_{1}^{*}[\Sht^{r}_{\cM_{\un d}}]=[\Sht^{r}_{\cM_{\un d,e}}]\in \Ch_{0}(\Sht^{r}_{\cM_{\un d,e}})
$.
\end{lemma}
\begin{proof}
This is because the maps $w_{1}$ and $h_{1}$ are both \'etale by Lemma \ref{l:uwh et}.
\end{proof}


\begin{lemma}
\source{higher-siegel-weil-unitary-nonsingular-terms}
\notready\label{l:u2}
\source{higher-siegel-weil-unitary-nonsingular-terms}
We have $u_{2}^{*}[\Sht^{r}_{\cM_{e}}]=[\Sht^{r}_{\cM_{\un d,e}}]\in \Ch_{0}(\Sht^{r}_{\cM_{\un d,e}})
$.
\end{lemma}
\begin{proof}
The diagram \eqref{rewrite ShtMde} is obtained from \eqref{rewrite ShtM} (for $\cM_{d}$ replaced by $\cM_{e}$) by base changing termwise along the map of the following two Cartesian diagrams induced by $u_{2}: \cA_{\un d,e}\to \cA_{e}$:
\begin{equation}
\xymatrix{\cA_{\un d,e}(k)\ar[r]\ar[d] & (\cA_{\un d,e})^{r+1}\ar[d]^{\Delta^{r+1}}="a" & \cA_{e}(k)\ar[r] \ar[d]_{}="b" & (\cA_{e})^{r+1}\ar[d]^{\Delta^{r+1}}\\
(\cA_{\un d,e})^{r+1}\ar[r]^{\Phi_{\cA_{\un d,e}}} & (\cA_{\un d,e})^{2r+2} & (\cA_{e})^{r+1}\ar[r]^{\Phi_{\cA_{e}}} & (\cA_{e})^{2r+2}
}
\end{equation}
Note that $u_{2}: \cA_{\un d,e}\to \cA_{e}$ is smooth since it exhibits $\cA_{\un d,e}$ as an open substack of a vector bundle over $\cA_{e}$ (using the condition $\mu_{\min}(\cE)>\max\{-d_{i}+3g'-2\}$). We conclude by applying Proposition \ref{p:base change} below.
\end{proof}


\subsubsection{Compatibility of cycle classes under Hitchin base change} To state the next result, we need some notations. Suppose we are given:
\begin{itemize}
\item stacks $S,M$ and $H$ that are locally of finite type over $k$ and can be stratified  into locally closed substacks that are global quotient stacks;
\item the stack $M$ is smooth of pure dimension $N$ with a map $f:M\to S$;
\item a map $\wt h: H  \to S^r \times_{\D^r, S^{2r}}M^{2r}$ (the fiber product uses the $r$-fold product of the diagonal  $\D^{r}: S^r \to S^{2r}$). 
\end{itemize}
Define $h \co H  \rightarrow M^{2r}$ as $\wt{h}$ followed by the second projection. Form the Cartesian square
\begin{equation}\label{ShtH}
\source{higher-siegel-weil-unitary-nonsingular-terms}
\xymatrix{\Sht_{H}\ar[d]\ar[r] &  H \times M \ar[d]^{h \times \Delta}\\
M^{r+1}\ar[r]^{\Phi_{M}} & M^{2r+2}}
\end{equation}
Let $u: S'\to S$ be a smooth representable morphism of pure relative dimension $D$. Let $M'=M\times_{S}S'$, $H'=H\times_{S^{r}}S'^{r}$ with natural maps $\wt h': H'\to S'^{r}\times_{\D^{r}, S'^{2r}}M'^{2r}$. Let $h':H'\to M'^{2r}$ be the resulting map. Let $u_{M}: M'\to M$ and $u_{H}: H'\to H$ be the natural maps. Form the Cartesian square
\begin{equation}\label{Sht'H}
\source{higher-siegel-weil-unitary-nonsingular-terms}
\xymatrix{\Sht'_{H}\ar[d]\ar[r] &  H'\ar[d]^{ h' \times \Delta} \times M' \\
M'^{r+1}\ar[r]^{\Phi_{M'}} & M'^{2r+2}}
\end{equation}
Since $S^{r+1}\times_{\Phi_{S}, S^{2r+2}, \D^{r+1}}S^{r+1}=S(k)$,  $\Sht_{H}$ decomposes as 
\begin{equation}
\Sht_{H}=\coprod_{s\in S(k)}\Sht_{H}(s).
\end{equation}
Similarly $\Sht'_{H}$ decomposes into the disjoint union of $\Sht_{H}'(s')$ indexed by $s'\in S'(k)$.
Then the natural map $u_{\Sht}: \Sht'_{H}\to \Sht_{H}$ lifts to an isomorphism
\begin{equation}
\Sht'_{H}=\Sht_{H}\times_{S(k)}S'(k)=\coprod_{s'\in S'(k)}\Sht_{H}(u(s')).
\end{equation}



\begin{prop}
\source{higher-siegel-weil-unitary-nonsingular-terms}
\notready\label{p:base change}  Let $\z\in \Ch_{*}(H)$ and $[M]$ be the fundamental class of $M$. Then we have
\source{higher-siegel-weil-unitary-nonsingular-terms}
\begin{equation}
u_{\Sht}^{*}\Phi_{M}^{!}(\z \times [M])=\Phi_{M'}^{!}u_{H}^{*}(\z \times [M])\in\Ch_{*-rN}(\Sht'_{H}).
\end{equation}
\end{prop}
%\mnote{Y: all references to Fulton should really be to Kresch. TF: adjusted the references}
\begin{proof}
Consider first the diagram where all squares are Cartesian
\begin{equation}
\xymatrix{\Sht_{H}\times_{S^{r}}S'^{r}\ar[r]\ar[d]^{v} & H' \times M' \ar[r]\ar[d]^{u_{H} \times u_M} & S'^{r+1}\ar[d]^{u^{r+1}}\\
\Sht_{H}\ar[r]\ar[d] & H \times M \ar[r]\ar[d]^{h \times \Delta} & S^{r+1}\\
M^{r+1}\ar[r]^-{\Phi_{M}} & M^{2r+2}
}
\end{equation}
Here the top vertical arrows are smooth and representable. By the compatibility of  Gysin map with flat pullback \cite[Theorem 2.1.12(ix)]{Kr99}, we have
\begin{equation}\label{vPhiu}
\source{higher-siegel-weil-unitary-nonsingular-terms}
v^{*}\Phi_{M}^{!}(\z \times [M])=\Phi_{M}^{!}u_{H}^{*}(\z \times [M])\in \Ch_{*-rN+(r+1)D}(\Sht_{H}\times_{S^{r+1}}S'^{r+1}).
\end{equation}
Here we recall that $D$ is the relative dimension of $u$.
%\mnote{Some equidimensionality assumption may be needed? TF: inserted this assumption} 
We have
\begin{equation}\label{ShtHSS'}
\source{higher-siegel-weil-unitary-nonsingular-terms}
\Sht_{H}\times_{S^{r+1}}S'^{r+1}=\coprod_{s\in S(k)}\Sht_{H}(s)\times (S'_{s})^{r+1}
\end{equation}
where $S'_{s}=u^{-1}(s)$, which is a smooth scheme over $k$. Factorize $u_{\Sht}$ as the composition
\begin{equation}
\begin{tikzcd}
\Sht'_{H} \ar[r, "{i}"] & \Sht_{H}\times_{S^{r}}(S'^{r}) \ar[r, "v"]  & \Sht_{H}.
\end{tikzcd}
\end{equation}
From \eqref{ShtHSS'} we see that $i$ is a regular embedding of codimension $(r+1)D$. Now $v$ and $u_{\Sht}$ are both smooth. Applying \cite[Theorem 2.1.12(ix)]{Kr99} we have $u^{*}_{\Sht}(-)=i^{!}v^{*}$ as maps $\Ch_{*}(\Sht_{H})\to \Ch_*(\Sht'_{H})$. Therefore
\begin{equation}\label{iuPhi}
\source{higher-siegel-weil-unitary-nonsingular-terms}
u_{\Sht}^{*}\Phi_{M}^{!}(\z \times [M])=i^{!}v^{*}\Phi_{M}^{!}(\z \times [M])\in \Ch_{*-rN}(\Sht_{H}').
\end{equation}	
On the other hand, consider the following diagram where all squares are Cartesian
\begin{equation}
\xymatrix{\Sht'_{H} \ar[r]^{i}\ar[d] & \Sht_{H}\times_{S^{r+1}}S'^{r+1}\ar[r]\ar[d] & H' \times M' \ar[d]^{ h' \times \Delta}\\
M'^{r+1}\ar[r]^-{\Phi_{1}} &  M^{r+1}\times_{\Phi_{S}\c f^{r+1}, S^{2r+2}}S'^{2r+2}\ar[r]^-{\Phi_{2}}\ar[d] & M'^{2r+2}\ar[d]^{u^{2r+2}_{M}}\\
& M^{r+1}\ar[r]^{\Phi_{M}} & M^{2r+2}
}
\end{equation}
Here $\Phi_{1}$ is the base change of $\Phi_{S'}$ and $\Phi_{2}$ is the base change of $\Phi_{M}$. The outer square of the top rows give \eqref{Sht'H}. By the transitivity of Gysin maps, we have
\begin{equation}
\Phi_{M'}^{!}u_{H}^{*}(\z\times [M])=\Phi^{!}_{1}\Phi^{!}_{2}u_{H}^{*}(\z \times [M]).
\end{equation}
Since $u^{2r+2}_{M}$  is smooth representable, we have $\Phi_{2}^{!}u_{H}^{*}(\z \times [M])=\Phi_{M}^{!}u_{H}^{*}(\z \times [M])$. Hence \begin{equation}
\Phi_{M'}^{!}u_{H}^{*}(\z\times [M])=\Phi^{!}_{1}\Phi_{M}^{!}u_{H}^{*}(\z\times [M]).
\end{equation}
Since both $\Phi_{1}$ and $i$ are regular embeddings of the same codimension, we have $\Phi^{!}_{1}(-)=i^{!}(-)$ as maps $\Ch_{*}(\Sht_{H}\times_{S^{r+1}}S'^{r+1})\to \Ch_{*-(r+1)D}(\Sht'_{H})$, by \cite[Theorem 2.1.12(xi)]{Kr99} and \cite[Theorem 6.2(c)]{Ful98}. Therefore
\begin{equation}\label{iPhiu}
\source{higher-siegel-weil-unitary-nonsingular-terms}
\Phi_{M'}^{!}u_{H}^{*}(\z\times [M])=i^{!}\Phi_{M}^{!}u_{H}^{*}(\z\times [M])\in \Ch_{*-rN}(\Sht'_{H}).
\end{equation}
Combining \eqref{vPhiu}, \eqref{iuPhi} and \eqref{iPhiu} we conclude
\begin{equation}
u_{\Sht}^{*}\Phi_{M}^{!}(\z\times [M])=i^{!}v^{*}\Phi_{M}^{!}(\z\times [M])=i^{!}\Phi_{M}^{!}u_{H}^{*}(\z\times [M])=\Phi_{M'}^{!}u_{H}^{*}(\z\times [M])\in \Ch_{*-rN}(\Sht'_{H}).
\end{equation}
\end{proof}


%
%---old -uses local product structure (may not be applicable to stacks?)----
%
%for every  $(\op\cL_{i}\xr{s}\cE,a)\in\cA_{\un d,e}(k)$ we can find an \'etale neighborhood $\cV$ and an \'etale map $\cU\to \cA_{e}\times V$ for some affine space $V$ over $k$. Base changing the diagram \eqref{rewrite ShtMde} over $\cU$, we get a diagram that is \'etale termwise over the termwise product of \eqref{rewrite ShtM} (with $d$ replaced by $e$) and the diagram
%\begin{equation}\label{Vr}
%\xymatrix{ V(k)\ar[d]\ar[r] & V^{r}\ar[d]^{\Delta^{r}}\\
%V^{r}\ar[r]^{\Phi_{V}} & V^{2r}}
%\end{equation}
%To prove the lemma, it suffices to prove the similar statement for this product diagram. 
%Write $M=\cM_{e}$, $S=\Sht^{r}_{\cM_{e}}$ and $H=\Hk^{1}_{\cM_{e}}$. Factorize the map $\Phi_{M\times V}: (M\times V)^{r}\to (M\times V)^{2r}$ into $M^{r}\times V^{r}\xr{\Phi_{M}\times\Id} M^{2r}\times V^{r}\xr{\Id\times\Phi_{V}}M^{2r}\times V^{2r}$, we get a diagram where all squares are Cartesian
%\begin{equation}
%\xymatrix{ & S\times V(k)\ar[d]\ar[rr]& & H^{r}\times V(k)\ar[d]\ar[rr] && H^{r}\times V^{r}\ar[d]^{((\pr_{0},\pr_{1})^{r}, \Delta^{r})}\\
%& M^{r}\times V^{r}\ar[rr]^-{\Phi_{M}\times\Id_{V^{r}}} && M^{2r}\times V^{r}\ar[rr]^-{\Id_{M^{2r}}\times\Phi_{V}} && M^{2r}\times V^{2r}}
%\end{equation}
%By the transitivity of the Gysin map, we have
%\begin{equation}\label{Phi1}
%\Phi_{M\times V}^{!}[H^{r}\times V^{r}]=(\Phi_{M}\times\Id_{V^{r}})^{!}(\Id_{M^{2r}}\times \Phi_{V})^{!}[H^{r}\times V^{r}]
%\end{equation}
%To compute $(\Id_{M^{2r}}\times \Phi_{V})^{!}[H^{r}\times V^{r}]$,  we may also view $H^{r}\times V(k)$ as obtained from $H^{r}\times V^{r}$ by base change along $\Phi_{V}: V^{r}\to V^{2r}$. By the compatibility of Gysin maps \cite[Theorem 6.2]{Fulton}, we see that
%\begin{equation}
% (\Id_{M^{2r}}\times \Phi_{V})^{!}[H^{r}\times V^{r}]= \pr_{V(k)}^{*}\Phi^{!}_{V}[V^{r}]
%\end{equation}
%Here $\pr_{V(k)}: H^{r}\times V(k)\to V(k)$ is the projection. Since \eqref{Vr} is a transverse intersection diagram, $\Phi^{!}_{V}[V^{r}]=[V(k)]$. We conclude that 
%\begin{equation}\label{Phi2}
% (\Id_{M^{2r}}\times \Phi_{V})^{!}[H^{r}\times V^{r}]=[H^{r}\times V(k)].
%\end{equation}
%For the same reason, viewing $S\times V(k)\to H^{r}\times V(k)$ as the base change along $\Phi_{M}: M^{r}\to M^{2r}$, and using the same compatibility of Gysin maps we get
%\begin{equation}\label{Phi3}
%(\Phi_{M}\times \Id_{V^{r}})^{!}[H^{r}\times V(k)]=\pr_{S}^{*}\Phi_{M}^{!}[H^{r}].
%\end{equation}
%Combing \eqref{Phi1}, \eqref{Phi2} and \eqref{Phi3} we get
%\begin{equation}
%\Phi_{M\times V}^{!}[H^{r}\times V^{r}]=\pr_{S}^{*}\Phi_{M}^{!}[H^{r}]
%\end{equation}
%which is what we wanted to prove (here $\pr_{S}$ is the transport of the map $u_{2}$ under the \'etale local chart $\Sht^{r}_{\cM_{\un d,e}}\leftarrow\Sht^{r}_{\cM_{\un d,e}}|_{\cU}\to \Sht^{r}_{\cM_{e}}\times V(k)$).



\begin{lemma}
\source{higher-siegel-weil-unitary-nonsingular-terms}
\notready\label{l:M und KR} Let  $(\{\cL_{i}\}_{1\le i\le n}, a')\in \cA_{\un d}(k)$ and $\cE':=\op_{i=1}^{n}\cL_{i}$.  Then we have an equality
\source{higher-siegel-weil-unitary-nonsingular-terms}
\begin{equation}
[\Sht^{r}_{\cM_{\un d}}]|_{\cZ^r_{\cE'}(a')}= \z^r_{\cL_{1},\cdots, \cL_{n}}(a')\in \Ch_{0}(\cZ^r_{\cE'}(a')).
\end{equation}
For the definition of $\z^r_{\cL_{1},\cdots, \cL_{n}}(a')$ see \S\ref{ssec: KR cycle classes}. 
\end{lemma}
\begin{proof}
We will apply the Octahedron Lemma \cite[Theorem A.10]{YZ} to a diagram of moduli stacks in our setting. Since the Octahedron Lemma requires certain stacks in question to be Deligne-Mumford, we need to rigidify our moduli stacks to satisfy these requirements. This is a minor technical issue which we encourage the reader to ignore: it is simply because the Octahedron Lemma in \cite{YZ} is not stated and proved in the most general form.




Let $v\in |X'|$. Let $P_{v}$ be the moduli space (a scheme!) of line bundles on $X'$ together with a trivialization of their fibers over $v$. Let $\G_{v}=\Res^{k_{v}}_{k}\G_{m}$. Then $P_{v}\to\Pic_{X'}$ is a $\G_{v}$-torsor.

Now for each moduli stacks $\cM_{\un d}, \cA_{\un d}$ and $\Hk^{r}_{\dot\cM_{\un d}}$ that involve an $n$-tuple of line bundles $\{\cL_{i}\}$, we write  $\dot\cM_{\un d}, \dot\cA_{\un d}$ and $\dot\Hk^{r}_{\dot\cM_{\un d}}$ to mean their rigidified versions where $\cL_{i}\in\Pic_{X'}$ is  replaced by $\dot\cL_{i}\in P_{v}$. Note that we do not impose any compatibility condition between the rigidifcation on $\cL_{i}$ and the rest of the structures classified by these moduli stacks. Define $\Sht^{r}_{\dot\cM_{\un d}}$ using the dotted version of the left one of the Cartesian diagrams in  \eqref{defn ShtMde}.

Note that $\dot\cM_{\un d}, \dot\cA_{\un d}$ and $\dot\Hk^{r}_{\dot\cM_{\un d}}$ are now schemes, and they are $\G_{v}^{n}$-torsors over their undotted counterparts. The dotted version of Lemma \ref{l:low deg M sm} remains valid if we add $n\deg(v)=\dim\G_{v}^{n}$ to the dimensions.  Also, $\Sht^{r}_{\dot\cM_{\un d}}\simeq \Sht^{r}_{\cM_{\un d}}\times_{\cA_{\un d}(k)}\dot\cA_{\un d}(k)$.  Since $\dot\cA_{\un d}(k)\to\cA_{\un d}(k)$ is surjective, to prove the Lemma, it suffices to prove its dotted version: for any $(\{\dot\cL_{i}\}, a')\in \dot\cA_{\un d}(k)$, writing $\cE':=\op_{i=1}^{n}\cL_{i}$, then there is an open and closed embedding $\cZ^r_{\cE'}(a')\incl \Sht^{r}_{\dot\cM_{\un d}}$ (using the rigifications $\dot\cL_{i}$ of $\cL_{i}$); then we shall prove
\begin{equation}\label{dot cycle eqn}
\source{higher-siegel-weil-unitary-nonsingular-terms}
[\Sht^{r}_{\dot\cM_{\un d}}]|_{\cZ^r_{\cE'}(a')}=\z^r_{\cL_{1},\cdots, \cL_{n}}(a')\in \Ch_{0}(\cZ^r_{\cE'}(a')).
\end{equation}



%The dotted analogues of Lemma \ref{l:uwh et}, \ref{l:ShtM isom} , \ref{l:u1} and  \ref{l:u2} remain true.


For $i=1,\cdots, n$,  let $\cN_{d_{i}}$ be the open substack of $\cM(1,n)$ consisting of points $(\cL\inj \cF,h)$ where $\chi(X',\cL)=-d_{i}$ and $\mu_{\min}(\cF)>-d_{i}+3g'-2$. Similarly define $\Hk^{1}_{\cN_{d_{i}}}$ and $\Sht^{r}_{\cN_{d_{i}}}$; these are open substacks of $\Hk^1_{\cM(1,n)}$ and $\Sht^r_{\cM(1,n)}$ respectively. 

Let $\dot\cN_{d_{i}}, \Hk^{1}_{\dot\cN_{d_{i}}}$ be the rigidified versions of $\cN_{d_{i}}$ and $\Hk^{1}_{\cN_{d_{i}}}$ where $\cL\in \Pic_{X'}$ is replaced by $\dot\cL\in P_{v}$. Let $\om_{i}: \dot\cN_{d_{i}}\to \Bun_{U(n)}$ and $\wt\om_{i}: \Hk^{1}_{\dot\cN_{d_{i}}}\to \Hk^{1}_{U(n)}$ be the forgetful maps.




We shall apply the Octahedron Lemma \cite[Theorem A.10]{YZ} to the following diagram:
\begin{equation}\label{9 term}
\source{higher-siegel-weil-unitary-nonsingular-terms}
\xymatrix{ (\Hk^{1}_{U(n)})^{r} \times \Bun_{U(n)}\ar[r]^-{\Delta}\ar[d]^{(\pr_{0},\pr_{1})^{r} \times \Delta} & \prod_{i=1}^{n} ((\Hk^{1}_{U(n)})^r \times \Bun_{U(n)} ) \ar[d]^{\prod ((\pr_{0},\pr_{1})^{r} \times \Delta)} &  \prod_{i=1}^{n} ((\Hk^{1}_{\dot\cN_{d_{i}}})^r \times \dot\cN_{d_{i}}) \ar[l]_-{\prod (\wt\om_{i}^{r} \times \omega_i)}\ar[d]^{\prod ((\pr_{0},\pr_{1})^{r} \times \Delta )}\\
\Bun_{U(n)}^{2r+2} \ar[r]^-{\Delta} & \prod_{i=1}^{n}\Bun_{U(n)}^{2r+2} & \prod_{i=1}^{n}\dot\cN_{d_{i}}^{2r+2}\ar[l]_-{\prod \om_{i}^{2r+2}}\\
\Bun_{U(n)}^{r+1} \ar[r]^-{\Delta}\ar[u]_{\Phi_{\Bun_{U(n)}}} & \prod_{i=1}^{n}\Bun_{U(n)}^{r+1}\ar[u]_{\prod\Phi_{\Bun_{U(n)}}}  & \prod_{i=1}^{n}\dot\cN_{d_{i}}^{r+1}\ar[u]_{\prod\Phi_{\dot\cN_{d_{i}}}}\ar[l]_-{\prod\om_{i}^{r+1}} 
}
\end{equation}
The fiber products of the three columns are
\begin{equation}\label{col Sht}
\source{higher-siegel-weil-unitary-nonsingular-terms}
\xymatrix{\Sht^{r}_{U(n)}\ar[r]^-{\Delta} & \prod_{i=1}^{n}\Sht^{r}_{U(n)}  & \prod_{i=1}^{n}\Sht^{r}_{\dot\cN_{d_{i}}}\ar[l]}
\end{equation}
where
\begin{equation}\label{ShtN decomp}
\source{higher-siegel-weil-unitary-nonsingular-terms}
\Sht^{r}_{\dot\cN_{d_{i}}}=\coprod_{\substack{\dot\cL_{i}\in P_{v}(k), \\ \chi(X',\cL_{i})=-d_{i}}}\left(\cZ^{r}_{\cL_{i}}(0)^{*}\coprod\left(\coprod_{ a'_{ii}\in \cA_{\cL_{i}}(k)}\cZ^{r}_{\cL_{i}}(a'_{ii})\right)\right).
\end{equation}

Let $\cMa_{\un d}$ be the moduli stack of $(\{\dot\cL_{i}\}, \op_{i=1}^{n}\cL_{i}\xr{t'} \cF,h)$ defined similarly as $\dot\cM_{\un d}$ but without the condition that $t'$ be injective, only that $t'_{i}=t|_{\cL_{i}}$ be injective. Then $\dot\cM_{\un d}\incl \cMa_{\un d}$ is open. Similarly define $\Hk^{1}_{\cMa_{\un d}}$ and $\Sht^{r}_{\cMa_{\un d}}$. Note that $\cMa_{\un d}$ is exactly the fiber product of
\begin{equation}\label{Mbar d}
\source{higher-siegel-weil-unitary-nonsingular-terms}
\xymatrix{\Bun_{U(n)}\ar[r]^-{\Delta} & \prod_{i=1}^{n}\Bun_{U(n)} & \prod_{i=1}^{n}\dot\cN_{d_{i}}\ar[l]_-{\prod\om_{i}}.}
\end{equation}
Similar remarks apply to $\Hk^{1}_{\cMa_{\un d}}$.  Therefore the fiber products of the three rows of \eqref{9 term} are
\begin{equation}\label{row Hk}
\source{higher-siegel-weil-unitary-nonsingular-terms}
\xymatrix{ (\Hk^{1}_{\cMa_{\un d}})^{r} \times \cMa_{\un d} \ar[d]^{(\pr_{0},\pr_{1})^{r} \times \Delta} \\
\cMa_{\un d}^{2r+2} \\
\cMa_{\un d}^{r+1}\ar[u]_{\Phi_{\cMa_{\un d}}}}
\end{equation}
The common fiber product of \eqref{col Sht} and \eqref{row Hk} is $\Sht^{r}_{\cMa_{\un d}}$, which decomposes as a disjoint union over the groupoid $\cB(k)$ of $(\dot\cL_{i},a'_{ii})_{1\le i\le n}$ where $\dot\cL_{i}\in P_{v}(k)$ with $\chi(X',\cL_{i})=-d_{i}$ and $a'_{ii}:\cL_{i}\to\s^{*}\cL_{i}^{\vee}$ injective Hermitian. For a point ${(\dot\cL_{i}, a'_{ii})_{1\le i\le n}}$ of $\cB(k)$, we have 
\begin{equation}
\Sht^{r}_{\cMa_{\un d}}|_{(\dot\cL_{i}, a'_{ii})_{1\le i\le n}} = \cZ^{r}_{\cL_{1},\cdots,\cL_{n}}(a'_{11},\cdots, a'_{nn})^{\circ},
\end{equation}
where $\Sht^{r}_{\cMa_{\un d}}|_{(\dot\cL_{i}, a'_{ii})_{1\le i\le n}}$ means the pullback of $
\Sht^{r}_{\cMa_{\un d}}$ to $\Spec k$ along the corresponding $\Spec k \rightarrow \cB(k)$. 

We check that the assumptions for applying the Octahedron Lemma are satisfied (the numbering below refers to that in \cite[Theorem A.10]{YZ}). 
\begin{enumerate}
\item All members in the diagram \eqref{9 term} are smooth and equidimensional. This is clear for $\Bun_{U(n)}$ and $\Hk^1_{U(n)}$. The same argument as in Lemma \ref{l:low deg M sm} proves that $\dot\cN_{d_{i}}$ and $\Hk^{1}_{\dot\cN_{d_{i}}}$ are smooth of pure dimension $d_{i}n+(n^2-2n+2)(g-1)+\deg(v)$.

\item We check that, in forming the fiber products of the middle and bottom rows and the left and middle columns,  the intersections are proper intersections with smooth equidimensional outcomes with the expected dimension. Here we use Lemma \ref{lem: smooth Bun_G} and \ref{lem: shtuka geometry} to argue for the left and middle columns. For the rows, the same argument as in Lemma \ref{l:low deg M sm} proves that $\cMa_{\un d}$ and $\Hk^{1}_{\cMa_{\un d}}$ are smooth of the same dimension as $\dot\cM_{\un d}$, which is $dn-(n^2-2n)(g-1)+n\deg(v)$. This is the virtual dimension for $\cMa_{\un d}$ as the fiber product of \eqref{Mbar d}, since 
\begin{align*}
\sum_{i=1}^{n}\dim \cN_{d_{i}}-(n-1)\dim \Bun_{U(n)}  &=\sum_{i=1}^{n} \left(d_{i}n+(n^{2}-2n+2)(g-1)+\deg(v)\right)-n^{2}(n-1)(g-1) \\
&=dn-(n^{2}-2n)(g-1)+n\deg(v).
\end{align*}

\item We check the fiber products of the top row and right column of \eqref{9 term} satisfy the conditions for  \cite[A.2.10]{YZ}. The fiber product of the top row is also a proper intersection: this follows from the same calculation as for the middle and bottom rows. The fiber product of the right column is also a proper intersection: this uses the decomposition \eqref{ShtN decomp} and the calculation of the dimension of $\cZ^{r}_{\cL_{i}}(0)^{*}$ in Proposition \ref{p:dim ZL0} and the dimension of $\cZ^{r}_{\cL_{i}}(a'_{ii})^{*}$ in Proposition \ref{p:ShM dim}.

The only issue is that $(\Hk^{1}_{\cMa_{\un d}})^{r+1}$ may not be a Deligne-Mumford stack, which was part of the requirement of \cite[A.2.10]{YZ}. However, we argue that this is not really an issue. The proof of the Octahedron Lemma allows the following flexibility: since eventually we only care about the $0$-cycles restricted to $\Sht^{r}_{\dot\cM_{\un d}}$, in the middle steps of forming the fiber products, we may restrict to open substacks as long as the final fiber product contains $\Sht^{r}_{\dot\cM_{\un d}}$ and only need check the relevant requirements there. Now in \eqref{row Hk} we may restrict to the open substack $(\Hk^{1}_{\dot\cM_{\un d}})^{r+1}\subset (\Hk^{1}_{\cMa_{\un d}})^{r+1}$, which is  a scheme.


\item The same remark as above shows that it suffices to check that the fiber squares obtained from \eqref{col Sht} and \eqref{row Hk}, after replacing $\cMa_{\un d}$ by $\dot\cM_{\un d}$,  each satisfy the condition \cite[A.2.8]{YZ}. Therefore it suffices to check
\begin{itemize}
\item $\Sht^{r}_{\dot\cM_{\un d}}$ admits a finite flat presentation in the sense of \cite[Definition A.1]{YZ}. This is true because $\Sht^{r}_{\dot\cM_{\un d}}$ is a scheme.
\item The diagonal map $\Delta: \Sht^{r}_{U(n)}\incl \prod_{i=1}^{n}\Sht^{r}_{U(n)}$ is a regular local immersion. This is true because $\Sht^{r}_{U(n)}$ is a smooth Deligne-Mumford stack.
\item The map $\Phi_{\dot\cM_{\un d}}: \dot\cM_{\un d}^{r+1}\to \dot\cM_{\un d}^{2r+2}$ is a regular local immersion. This is true because $\dot\cM_{\un d}$ is a smooth equidimensional scheme by the dotted version of Lemma \ref{l:low deg M sm}. %Since we only care about $\Sht^{r}_{\cM_{\un d}}$ in the end, we only need to check the conditions of \cite[A.2.8]{YZ} for the diagram obtained from \eqref{row Hk} by replacing $\cMa_{\un d}$ with $\cM_{\un d}$. The desired property for $\cM_{\un d}$ follows from the fact that it is a smooth Deligne-Mumford stack according to Lemma \ref{l:low deg M sm}. \mnote{TF: completed this bullet point}
\end{itemize}
\end{enumerate}
The conclusion of the (variant of) Octahedron Lemma says that the following two elements in $\Ch_{0}(\Sht^{r}_{\cMa_{\un d}})$
\begin{equation}
\Phi_{\cMa_{\un d}}^{!}\Delta_{(\Hk^{1}_{U(n)})^{r}}^{!}[\prod_{i=1}^{n}(\Hk^{1}_{\dot\cN_{d_{i}}})^{r} \times \dot\cN_{d_{i}}] \quad \mbox{and}\quad \Delta_{\Sht^{r}_{U(n)}}^{!}(\prod\Phi_{\dot\cN_{d_{i}}})^{!}[\prod_{i=1}^{n}(\Hk^{1}_{\dot\cN_{d_{i}}})^{r} \times \dot\cN_{d_{i}}]
\end{equation}
become the same when restricted to $\Sht^{r}_{\dot\cM_{\un d}}$. Further restricting to $\cZ_{\cE'}^r(a')$ we get the desired identity \eqref{dot cycle eqn}. %\mnote{TF: added the superscript $r$ to $\cZ$ in many places}
\end{proof}



\begin{proof}[Proof of Theorem \ref{th:compare KR weak}] Restricting the equality in Lemma \ref{l:M und KR} to $\cZ_{\cE}^r(a)$, which is open and closed in $\cZ_{\cE'}^r(a')$ by Corollary \ref{c:ZE open closed}, we get
\begin{equation}\label{ZL}
\source{higher-siegel-weil-unitary-nonsingular-terms}
\z^r_{\cL_{1},\cdots, \cL_{n}}(a')|_{\cZ^r_{\cE}(a)}=[\Sht^{r}_{\cM_{\un d}}]|_{\cZ_{\cE}^r(a)}.
\end{equation}
For fixed $(\op\cL_{i}\inj \cE,a)\in \cA_{\un d,e}(k)$, $\cZ^{r}_{\cE}(a)$ can be viewed as a finite \'{e}tale cover of an open-closed substack in $\Sht^{r}_{\cM_{\un d}},\Sht^{r}_{\cM_{\un d,e}}$ and $\Sht^{r}_{\cM_{e}}$ by Lemma \ref{l:ShtM isom}. By Lemma \ref{l:u1} and Lemma \ref{l:u2} we have
\begin{equation}
[\Sht^{r}_{\cM_{\un d}}]|_{\cZ^r_{\cE}(a)}=[\Sht^{r}_{\cM_{\un d,e}}]|_{\cZ^r_{\cE}(a)}=[\Sht^{r}_{\cM_{e}}]|_{\cZ^{r}_{\cE}(a)}.
\end{equation}
Combining this with \eqref{ZL} proves the theorem.
\end{proof}


\subsection{Proof of Theorem \ref{th:compare KR} for $X'=X\coprod X$ or $X_{k'}$}\label{ss:compare cycle split} Here $k'/k$ is the quadratic extension. 
\source{higher-siegel-weil-unitary-nonsingular-terms}

In the case $X'=X\coprod X$,  we have $\Bun_{U(n)}\cong \Bun_{\GL_{n}}$. We shall identify a Hermitian bundle $\cF$ on $X'$ with a pair of vector bundles $(\cF_{1},\cF_{2})$ equipped with an isomorphism $\cF_{2}\cong \cF_{1}^{\vee}$, each living on one copy of $X$.  A vector bundle $\cE$ on $X'$ of rank $n$ corresponds to two rank $n$ vector bundles $(\cE_{1},\cE_{2})$, each living on one copy of $X$. Now $\cA_{\cE}(k)$ is the set of injective maps $a: \cE_{1}\to \cE^{\vee}_{2}$.  A good framing $s:\op_{i=1}^{n}\cL_{i}\inj \cE$ for $(\cE,a)$ now consists of line bundles $\cL_{i}=(\cL_{i,1}, \cL_{i,2})$ ($1\le i\le n$) satisfying the same conditions in Definition \ref{def: auxiliary E'}; it is called {\em very good} if it satisfies the additional conditions
\begin{enumerate}
\item[($3_{1}$)] $\mu_{\min}(\cE_{1})>\max\{\deg\cL_{i,1}+2g-1\}_{1\le i\le n}$, and
\item[($3_{2}$)] $\mu_{\min}(\cE_{2})>\max\{\deg\cL_{i,2}+2g-1\}_{1\le i\le n}$.
\end{enumerate}
The same argument of Lemma \ref{l:weak implies strong} shows that it suffices to prove the analogue of Theorem \ref{th:compare KR weak}, i.e., prove Theorem \ref{th:compare KR} for very good framings.

In both the $X'=X\coprod X$ and $X'=X_{k'}$ case, we need to modify the definitions of $\cM_{\un d}$ and $\cM_{\un d,e}$ as follows. In the definition of $\cM_{\un d,e}$, we use the notion of very good framing just defined over geometric fibers of $X'_{S}\to S$ (which are of the form $X_{\ov s}\coprod X_{\ov s}$).  In the definition of $\cM_{\un d}$, we change the inequality \eqref{slope F} to two inequalities over the geometric fibers of $X'_{S}\to S$
\begin{eqnarray*}
\mu_{\min}(\cF_{1})>\max\{\deg\cL_{i,1}+2g-1\}_{1\le i\le n},\\
\mu_{\min}(\cF_{2})>\max\{\deg\cL_{i,2}+2g-1\}_{1\le i\le n}.
\end{eqnarray*}
The same inequalities should be imposed in the definition of $\cN_{d_{i}}$ that appear in the proof of Lemma \ref{l:M und KR}. With these changes, the argument for proving Theorem \ref{th:compare KR weak} goes through.
